\section{Related Work}

\subsection{RLHF Evaluation Methods}

Reinforcement Learning from Human Feedback (RLHF) has become the dominant paradigm for aligning language models to human preferences. InstructGPT \cite{ouyang2022instructgpt} introduced the standard three-stage pipeline: supervised fine-tuning (SFT), reward model training from pairwise comparisons, and policy optimization via PPO. Evaluation focused on preference win rates: InstructGPT achieved 85\% preference agreement vs GPT-3 base model on human evaluations. Constitutional AI \cite{bai2022constitutional} extended this with AI feedback for harmlessness, reporting improved safety scores while maintaining helpfulness. Both methods treat \textbf{preference convergence as success} — higher win rates indicate better alignment.

However, these benchmarks measure only \textbf{unidirectional alignment} (AI$\rightarrow$human): does the model produce outputs humans prefer? They do not measure \textbf{bidirectional alignment} (human$\rightarrow$AI): do humans preserve critical evaluation capacity when exposed to aligned models? High preference agreement could indicate quality improvement (users genuinely prefer better responses) or homogenization (users habituate to model style and stop critically evaluating). Existing metrics cannot distinguish these scenarios. \textbf{Preference convergence is appropriate for objective tasks} (math, factual QA with clear ground truth) where agreement indicates learning correct answers. Our concern applies to \textbf{subjective tasks} (creative writing, opinion questions, stylistic preferences) where diverse preferences are legitimate and convergence may signal homogenization.

\textbf{Literature Search Methodology:} We searched ACL Anthology, arXiv (cs.CL, cs.LG), and Google Scholar for combinations of ``RLHF diversity'', ``preference variance'', ``preference entropy'', ``annotator agreement entropy'', ``human feedback diversity'', and ``bidirectional alignment'' (2020-2026). We found work measuring inter-annotator agreement (Cohen's kappa, Fleiss' kappa) and variance in recommender systems but no prior work applying Shannon entropy to RLHF preference distributions as a bidirectional alignment metric.

\textbf{WebGPT} \cite{nakano2021webgpt} and \textbf{OpenAI Summarization} \cite{stiennon2020summarize} provide relevant precedents. OpenAI Summarization collected multi-annotator ratings on the same summary candidates, enabling variance analysis. However, their published metrics still focused on mean preference scores, not entropy or diversity measures. This suggests the data structure for diversity measurement exists in some datasets (multi-annotator fixed responses) but has not been leveraged for bidirectional alignment evaluation.

\subsection{Bidirectional Alignment and Human Agency}

Human-AI interaction research emphasizes user agency and empowerment \cite{amershi2019guidelines}. The ``Guidelines for Human-AI Interaction'' framework includes principles like ``support efficient correction'' and ``encourage granular feedback'' — both requiring that users maintain critical evaluation of AI outputs. However, these guidelines rely on \textbf{self-reported measures} (user surveys, perceived control) or \textbf{interaction patterns} (override rates, feedback frequency). Our entropy approach provides a \textbf{behavioral proxy} computable from existing preference data without new human evaluation.

Reward hacking literature \cite{skalse2022misalignment} addresses over-optimization on proxy metrics where models exploit misspecified reward functions. Our concern is orthogonal: not model behavior pathology (exploiting rewards) but \textbf{user behavior homogenization} (converging preferences). Entropy collapse would indicate alignment succeeded at making users agree, but failed at preserving legitimate diversity on subjective tasks.

\subsection{Information-Theoretic Diversity Measures}

Shannon entropy \cite{shannon1948information} H = -$\Sigma$ $p_i$ log($p_i$) is the foundational measure of distribution uncertainty. Higher entropy indicates greater diversity/unpredictability; lower entropy indicates concentration/consensus. Entropy has been applied to measure diversity in ecological systems (species distribution), information retrieval (term frequency), and machine learning (policy exploration). To our knowledge, its application to \textbf{human preference distributions} in RLHF alignment evaluation has not been explored in prior work (see Section 2.1 literature search).

Prior work on \textbf{preference diversity} in recommender systems \cite{nguyen2014exploring} showed that personalized algorithms can create filter bubbles by reducing content diversity. This parallels our concern: RLHF may reduce preference diversity by training users to converge on model-preferred responses. The key difference is \textbf{measurement level}: recommender systems track individual user exposure diversity, while we propose \textbf{population-level preference entropy} as collective critical evaluation proxy.

\subsection{Dataset Format and Annotation Cost Trade-offs}

Pairwise comparison collection (Bradley-Terry models, Thurstone scaling) is standard in preference elicitation due to \textbf{annotation efficiency}: labelers compare two options (binary choice) rather than rating multiple candidates independently. Anthropic-HH \cite{bai2022constitutional} and WebGPT \cite{nakano2021webgpt} use this format: each example presents one chosen and one rejected response for a given prompt. This minimizes labeler cognitive load and enables large-scale collection (160K+ comparisons for Anthropic-HH).

However, pairwise formats create a \textbf{diversity measurement incompatibility}. To compute per-prompt preference entropy, we need \textbf{multiple annotators rating the same response candidates} — a distribution over fixed options (e.g., ``40\% prefer A, 35\% prefer B, 25\% prefer C'' $\rightarrow$ H $\approx$ 1.05 nats). Anthropic-HH instead provides \textbf{unique response pairs per comparison} (A1 vs B1, A2 vs B2, ...), preventing aggregation into distributions.

\textbf{OpenAI Summarization} \cite{stiennon2020summarize} provides a counter-example: multiple annotators rated the same summary candidates (4-5 summaries per article, 3-5 annotators per summary). This multi-annotator fixed-response format enables entropy computation but incurs higher annotation cost (N annotators $\times$ M responses per prompt vs 1 annotator per pairwise comparison). Our contribution is documenting this \textbf{format-measurement trade-off} in the context of RLHF benchmarks and proposing alternatives when multi-annotator data is unavailable.

\subsection{Positioning Our Work}

We introduce \textbf{preference entropy as bidirectional alignment metric}, distinguishing our work from:
\begin{itemize}
\item \textbf{RLHF evaluation} (InstructGPT, Constitutional AI): Measures preference agreement (unidirectional AI$\rightarrow$human), not diversity (bidirectional human agency preservation).
\item \textbf{HCI user empowerment metrics}: Requires self-reported surveys; entropy uses behavioral data from existing benchmarks.
\item \textbf{Reward hacking detection}: Focuses on model over-optimization; we address user homogenization.
\end{itemize}

Our \textbf{dataset structure documentation} (pairwise-unique vs multi-annotator-fixed) is a methodological contribution: documenting data format requirements for diversity measurement in RLHF evaluation, tested on Anthropic-HH (n=100 prompts) with format-based inference for other datasets. This explains why existing benchmarks cannot retrospectively measure entropy without structural changes or alternative proxies.
